\section{Related Work}

\noindent\textbf{3D Reconstruction.} 
Recent 3D reconstruction research can be broadly categorized into traditional geometry-based and deep learning methods. The former relies on multi-view stereo (MVS)~\cite{YAN} and structure from motion (SfM)~\cite{sfm} to estimate scene depth and camera poses, producing point clouds and subsequent surface meshes. The latter integrates implicit functions (e.g., SDF, Occupancy)~\cite{Tri} with volumetric rendering for high-fidelity reconstruction, as exemplified by Neural Radiance Fields (NeRF)~\cite{nerf}. However, NeRF-based approaches often struggle with real-time performance in large-scale or dynamic scenarios. In contrast, 3D Gaussian Splatting~\cite{3dgs} encodes scenes as 3D Gaussians (with position, scale, and color), using differentiable point-based rendering to achieve fast training and inference while balancing accuracy and quality. Balancing high fidelity, scalability, and real-time capability remains a key challenge in 3D reconstruction.
Within the field of 3D reconstruction, there are primarily two main sub-tasks: novel view synthesis (NVS) and surface reconstruction. 

\vspace{0.5em}

\noindent\textbf{Novel View Synthesis.} 
Novel View Synthesis (NVS) aims to generate a target image from an arbitrary camera pose, given source images and their camera poses~\cite{light, lum}. NeRF~\cite{nerf} integrates implicit representations with volume rendering\cite{volume, eff}, demonstrating impressive results in view synthesis. However, dense point sampling remains a major bottleneck for rendering speed. To address this, various methods accelerate NeRF by replacing the original multi-layer perceptrons (MLPs)~\cite{learining, deep} with discretized representations, such as voxel grids~\cite{dir}, hash encodings~\cite{Instant}, or tensor radiation fields~\cite{tensorf}. Additionally, some approaches~\cite{bake, merf} distill pretrained NeRFs into sparse representations, enabling real-time rendering. Recent advancements in 3D Gaussian Splatting (3DGS) have significantly improved real-time rendering, demonstrating that continuous representations are not strictly necessary. However, directly optimizing and rendering at high resolutions drastically increase memory overhead, making it challenging to achieve real-time reconstruction of high-quality scenes on mainstream GPUs with limited memory (24GB). Our approach specifically addresses this challenge by reducing the computational cost of high-resolution processing while preserving reconstruction fidelity. 

\vspace{0.5em}

\textbf{Multi-View Surface Reconstruction.} 
In recent years, multi-view reconstruction methods have evolved from traditional geometric approaches to neural implicit representations. Traditional multi-view stereo methods~\cite{re_MVS1,re_MVS2,re_MVS3,re_MVS4,re_MVS5,re_MVS6} primarily rely on extracting dense depth maps~\cite{re_depth1,re_depth2} from multiple images, fusing them into point clouds~\cite{re_point,re_point2}, and generating scene models through triangulation or implicit surface fitting~\cite{re_fit}. Although these techniques have been widely adopted in both academia and industry, they are often susceptible to artifacts and loss of detail due to matching~\cite{re_matching1} errors, noise, and local optimum issues during reconstruction. 
Recent advances in neural implicit representations, such as NeRF~\cite{re_nerf} and its SDF-based variants~\cite{NeuS,MonoSDF}, have shifted the reconstruction paradigm from explicit depth estimation toward learning continuous volumetric or surface fields directly from images. These approaches inherently model geometry and appearance jointly, offering better robustness to occlusions and textureless regions.
While neural implicit methods offer high-quality reconstructions, they remain computationally intensive and face scalability challenges. As an alternative, 3D Gaussian Splatting (3DGS)~\cite{3dgs} adopts an explicit representation by projecting anisotropic Gaussians onto image space, enabling efficient and differentiable rasterization~\cite{yifan2019differentiable}. Despite its ability to support real-time rendering with high visual fidelity, 3DGS often suffers from insufficient geometric supervision, particularly in sparse-view or large-scale scenarios~\cite{neusg}. To address this limitation, recent methods such as VCR-GauS~\cite{vcr-gaus}, Vastgaussian~\cite{vast}, and SuGaR~\cite{SuGaR} introduce view-consistent depth and normal constraints~\cite{turkulainen2024dnsplatter,bae2024dsine,yin2019enforcing}, significantly enhancing both reconstruction accuracy and convergence stability. These developments position 3DGS as a promising solution for high-resolution and scalable surface reconstruction.

%With the advancement of deep learning, NeRF~\cite{re_nerf}and its variants have emerged, modeling the color and density in continuous space using multi-layer perceptrons and generating novel views via volume rendering, which significantly enhances rendering quality.Subsequent methods aimed to optimize both quality and speed. Techniques like position encoding and band-limited coordinate networks are combined with neural radiance fields for pre-filtered scene representation~\cite{re_mip,re_bacon}. Innovations to speed up rendering included leveraging spatial data structures and adjusting MLP size~\cite{re_mlpsize1,re_mlpsize2,re_mlpsize3,re_mlpsize4,re_mlpsize5}. However, high-quality rendering with NeRF also comes at the cost of increased computational overhead and longer training times, posing challenges for real-time applications and large-scale scenarios~\cite{Tri-miprf,mega-nerf}. To mitigate these limitations, subsequent research has proposed various acceleration strategies, such as Zip-NeRF~\cite{zip-nrf}, InstantNGP~\cite{InstantNGP}, and Plenoxels~\cite{plenoxels}, which employ techniques like low-pass filtering, multi-resolution hash grids, and sparse voxel grids to substantially boost computational efficiency~\cite{grid-nerf} while preserving rendering quality.


